\documentclass[11pt]{article}
\usepackage[T1]{fontenc}
\usepackage{lmodern}
\usepackage[margin=1in]{geometry}
\usepackage{amsmath,amssymb,amsthm}
\usepackage{booktabs,array}
\usepackage[hidelinks]{hyperref}

\newcommand{\llbracket}{[\![}
\newcommand{\rrbracket}{]\!]}

\hypersetup{
  pdftitle={OFF-Cell Counts in Wolfram Rule 11: Proof and Lean 4 Formalization of Colin Barker's OEIS A266258 Conjectures},
  pdfauthor={Derek Blair}
}

\newtheorem{theorem}{Theorem}[section]
\newtheorem{proposition}[theorem]{Proposition}
\newtheorem{lemma}[theorem]{Lemma}
\newtheorem{corollary}[theorem]{Corollary}
\theoremstyle{definition}
\newtheorem{definition}[theorem]{Definition}
\newtheorem{remark}[theorem]{Remark}

\newcommand{\N}{\mathbb{N}}
\newcommand{\Z}{\mathbb{Z}}
\newcommand{\Q}{\mathbb{Q}}
\newcommand{\seq}[1]{\texttt{#1}}
\newcommand{\ruleTT}{\operatorname{rule}_{11}}
\newcommand{\ca}{\operatorname{ca}}

\setlength{\emergencystretch}{2em}

\title{OFF-Cell Counts in Wolfram Rule 11:\\[2pt]
\large Proof and Lean~4 Formalization of Colin Barker's OEIS \seq{A266258} Conjectures}
\author{Derek Blair}
\date{24 November 2026}

\begin{document}
\maketitle

\begin{abstract}
Let $\ca : \N \times \Z \to \{\mathtt{false}, \mathtt{true}\}$ denote the bi-infinite evolution of Wolfram's Elementary Cellular Automaton Rule~11 ($11 = 00001011_2$) starting from a single ON cell at the origin, $\ca(0, j) = \mathrm{decide}(j = 0)$, and let $a(n) = \seq{A266258}(n)$ count the OFF (white) cells inside the active light-cone window $[-n, n]$ at generation $n \ge 0$. In December~2015 and April~2019, Colin Barker posted a three-part conjecture bundle on OEIS~\seq{A266258}: the rational closed form $a(n) = \bigl(2(-1)^n n + 2n - 3(-1)^n + 1\bigr)/2$ for all $n > 0$, the four-step linear recurrence $a(n) = 2a(n-2) - a(n-4)$ for all $n > 4$, and the rational ordinary generating function $\sum_{n \ge 0} a(n)\,X^n = X(2 + 3X - 2X^2 + X^3)/\bigl((1 - X)^2(1 + X)^2\bigr)$.

We prove Barker's full conjecture bundle unconditionally together with the explicit initial value $a(0) = 0$, both the cleared-denominator and explicit unit-inverse generating-function identities in $\Q\llbracket X \rrbracket$, and exact sharpness of both domain guards ($n = 0$ and $n = 4$). Unlike defect-annihilating rules such as Rule~7 (\seq{A266220}/\seq{A266222}) and Rule~23 (\seq{A266437}/\seq{A266439}), Rule~11 breaks left-right symmetry via the Boolean normal form $\ruleTT(l, c, r) = (\neg l) \land (\neg c \lor r)$ and sustains a \emph{persistent two-cell glider defect} at $\{n - 1, n\}$ moving rightward at speed~$1$ against a period-$2$ alternating exterior on all of $\Z$: for every $n \ge 1$ and $j \in \Z$, $\ca(n, j) = \operatorname{xor}\bigl(\mathrm{decide}(\mathrm{Odd}(n)),\, \mathrm{decide}(j \in \{n - 1, n\})\bigr)$. Restricting the OFF-cell locus to $[-n, n]$ yields the two-element OFF pair $\{2m, 2m + 1\}$ ($a(2m + 1) = 2$ for $m \ge 0$) and the contiguous OFF interval $[-2m, 2m - 2]$ ($a(2m) = 4m - 1$ for $m \ge 1$), from which the rational closed form, four-step recurrence, and cleared- and unit-inverse generating-function identities in $\Q\llbracket X \rrbracket$ follow directly. We also prove the exact OFF/ON complement law $a(n) + a_{\mathrm{ON}}(n) = 2n + 1$, establish that the isolated-zero pattern $(1, 0, 1)$ never occurs in $\ca$ (yielding bit-for-bit space-time identity between Rules~11 and~43), contrast Rule~11 with Rule~47 (\seq{A266664}), construct an $O(n^2)$ causal-cone mirror confirming $a(0), \dots, a(12)$ by pure kernel reduction, and formalize both light-cone and fixed-box zero-clamping traps. Every theorem is machine-checked in Lean~4 (\texttt{v4.34.1}, Mathlib \texttt{d13f23b7}) with zero \texttt{sorry}s, allowlist-only axioms, independent \texttt{leanchecker} kernel replay, and zero computational trust escapes.
\end{abstract}

\section{Introduction, Primary Sources, and Convention Audit}

\subsection{Wolfram Rule 11 and OEIS \texorpdfstring{\seq{A266258}}{A266258}}

In December~2015, Robert Price cataloged the single-seed evolution of Wolfram's Elementary Cellular Automaton (ECA) Rule~11~\cite{wolfram2002,mathworldECA} across seven companion OEIS sequences \seq{A266253}--\seq{A266259}. Specifically, \seq{A266253}~\cite{oeis266253} lists the triangle of cell states in the active light-cone window $[-n, n]$, \seq{A266254} and \seq{A266255}~\cite{oeis266254,oeis266255} encode the row bitstrings in base~$10$ and base~$2$, \seq{A266256}~\cite{oeis266256} counts ON (black) cells, \seq{A266257} and \seq{A266259} track cumulative ON and OFF totals, and \seq{A266258}~\cite{oeis266258} counts OFF (white) cells in row~$n$:
\[
a(n) = 0, 2, 3, 2, 7, 2, 11, 2, 15, 2, 19, 2, 23, 2, 27, 2, 31, 2, 35, 2, \dots \qquad (n \ge 0).
\]

Under Wolfram's standard binary rule numbering~\cite{wolfram2002}, an ECA rule code $c \in \{0, \dots, 255\}$ maps a binary neighborhood $(l, c_0, r) \in \{0, 1\}^3$ ($1 = \text{ON/black}$, $0 = \text{OFF/white}$) to bit $4l + 2c_0 + r$ of~$c$. For Rule~11,
\[
11 = 8 + 2 + 1 = 00001011_2,
\]
whose explicit $8$-entry truth table appears in Table~\ref{tab:rule11-tt}.

\begin{table}[ht]
\centering
\caption{Wolfram Elementary Cellular Automaton Rule~11 ($11 = 00001011_2$) truth table. Output is~$1$ if and only if $l = 0$ and $(c = 0 \lor r = 1)$, i.e., $\ruleTT(l, c, r) = (\neg l) \land (\neg c \lor r)$.}
\label{tab:rule11-tt}
\smallskip
{\small
\begin{tabular}{@{}c c c c@{}}
\toprule
Neighborhood $(l, c, r)$ & Binary index $4l + 2c + r$ & Bit of $11$ & $\ruleTT(l, c, r)$ \\
\midrule
$(1, 1, 1)$ & $7$ & $0$ & $0$ \\
$(1, 1, 0)$ & $6$ & $0$ & $0$ \\
$(1, 0, 1)$ & $5$ & $0$ & $0$ \\
$(1, 0, 0)$ & $4$ & $0$ & $0$ \\
$(0, 1, 1)$ & $3$ & $1$ & $1$ \\
$(0, 1, 0)$ & $2$ & $0$ & $0$ \\
$(0, 0, 1)$ & $1$ & $1$ & $1$ \\
$(0, 0, 0)$ & $0$ & $1$ & $1$ \\
\bottomrule
\end{tabular}
}
\end{table}

In the \texttt{\%F} section of OEIS~\seq{A266258} (live revision~\#14, February~16, 2025~\cite{oeis266258}), Colin Barker posted the following three-part conjecture bundle on December~27, 2015 and April~14, 2019:
\begin{quote}\itshape
Conjectures from Colin Barker, Dec 27 2015 and Apr 14 2019:\\
$a(n) = (2(-1)^n n + 2n - 3(-1)^n + 1)/2$ for $n > 0$\\
$a(n) = 2a(n - 2) - a(n - 4)$ for $n > 4$.\\
G.f.: $x(2 + 3x - 2x^2 + x^3) / ((1 - x)^2(1 + x)^2)$.
\end{quote}
Table~\ref{tab:small-n} summarizes generations $n = 0, \dots, 6$, illustrating the period-$2$ alternating exterior, the right-edge two-cell defect inside $[-n, n]$, the OFF-cell counts $a(n)$, Barker's three formulas, and the complementary ON-cell counts $a_{\mathrm{ON}}(n) = \seq{A266256}(n)$.

\begin{table}[ht]
\centering
\caption{Generations $n = 0, \dots, 6$ of Rule~11 evolved on all of $\Z$ from the single ON cell at $j = 0$, alongside Barker's closed form, four-step recurrence, generating function coefficients, and companion ON-cell count $a_{\mathrm{ON}}(n) = \seq{A266256}(n)$.}
\label{tab:small-n}
\smallskip
{\small
\begin{tabular}{@{}r c l r r c r r@{}}
\toprule
$n$ & Exterior & Window $[-n, n]$ & $a(n)$ & Closed form & $2a(n-2) - a(n-4)$ & $[X^n]F(X)$ & $a_{\mathrm{ON}}(n)$ \\
\midrule
$0$ & $0$ & \texttt{1} & $0$ & $-1^*$ & --- & $0$ & $1$ \\
$1$ & $1$ & \texttt{100} & $2$ & $2$ & --- & $2$ & $1$ \\
$2$ & $0$ & \texttt{00011} & $3$ & $3$ & --- & $3$ & $2$ \\
$3$ & $1$ & \texttt{1111100} & $2$ & $2$ & --- & $2$ & $5$ \\
$4$ & $0$ & \texttt{000000011} & $7$ & $7$ & $6^\dagger$ & $7$ & $2$ \\
$5$ & $1$ & \texttt{11111111100} & $2$ & $2$ & $2$ & $2$ & $9$ \\
$6$ & $0$ & \texttt{0000000000011} & $11$ & $11$ & $11$ & $11$ & $2$ \\
\bottomrule
\end{tabular}

\smallskip
{\footnotesize ${}^*$\,$n = 0$ is excluded by $n > 0$.\quad ${}^\dagger$\,$n = 4$ gives $2(3) - 0 = 6 \ne 7$ and is excluded by $n > 4$.}
}
\end{table}

\subsection{Convention Audit and Formalization Traps}
\label{sec:conventions}

A first-hand audit of Robert Price's Mathematica program (\texttt{\%t}), offset (\texttt{\%O A266258 0,2}), and b-file \texttt{b266258.txt} ($n = 0, \dots, 999$)~\cite{oeis266258} fixes five conventions that must be preserved verbatim in any faithful specification:
\begin{enumerate}
\item \textbf{Evolve-on-$\Z$-first semantics:} Price's Mathematica code invokes {\small\nolinkurl{CellularAutomaton[11,{{1},0},rows-1,{All,All}]}} on the bi-infinite lattice before slicing {\small\nolinkurl{Take[ca[[k]],{rows-k+1,rows+k-1}]}} (width $2n + 1$ at $k = n + 1$) and taking {\small\nolinkurl{Length[catri[[k]]]-Total[catri[[k]]]}}. Because $\ruleTT(0,0,0) = 1$ and $\ruleTT(1,1,1) = 0$, the uniform exterior $|j| > n$ flips $0 \mapsto 1 \mapsto 0 \mapsto \cdots$ on every step (in particular, counting OFF cells on all of $\Z$ would be infinite at every even $n$, so truncation to $[-n, n]$ after $\Z$-evolution is essential). Both natural zero-clamping shortcuts fail immediately:
\begin{itemize}
\item \emph{Trap~A (light-cone zero-clamping):} clamping $|j| > m$ to $0$ between steps feeds $(0,0,1) \mapsto 1$ into $j = -2$ at step $n = 2$, turning row~$2$ from \texttt{00011} ($a(2) = 3$) into \texttt{10011} (OFF count~$2 \ne 3$).
\item \emph{Trap~B (fixed-box $[-n, n]$ zero-clamping):} evolving on the fixed box $[-n, n]$ with zero exterior corrupts row~$2$ from \texttt{00011} to \texttt{10010} (spurious ON bit at $j = -2$ cancelling lost ON bit at $j = +2$, accidentally preserving OFF count~$3$) and yields wrong OFF count $4 \ne 2 = a(3)$ at $n = 3$ (\texttt{0011100} vs.\ \texttt{1111100}).
\end{itemize}
Both traps are formalized in Proposition~\ref{prop:clamping-traps}. Hence $\ca : \N \to \Z \to \mathsf{Bool}$ must evolve on all of $\Z$ first and restrict to $\mathrm{Icc}(-n, n)$ only inside $a(n)$.
\item \textbf{Wolfram bit ordering and chirality:} Bit $4l + 2c + r$ of $11 = 00001011_2$ is confirmed against \seq{A266253} row~$1 = \mathtt{100}$ and row~$2 = \mathtt{00011}$ (reversing to $l + 2c + 4r$ gives mirror Rule~81, \seq{A266982}, with row~$1 = \mathtt{001}$). The explicit $8$-branch truth table $\ruleTT$ is pinned verbatim and proved equal to $\mathrm{Nat.testBit}\;11\;(4l + 2c + r)$.
\item \textbf{Exact guard boundaries:} Offset is $0$ with $a(0) = 0$ pinned explicitly. The guard $n > 0$ in the closed form is sharp because at $n = 0$ Barker's formula evaluates to $(0 + 0 - 3 + 1)/2 = -1 \ne 0$. The guard $n > 4$ in the four-step recurrence is sharp because at $n = 4$ we have $2a(2) - a(0) = 2(3) - 0 = 6 \ne 7 = a(4)$ (creating the $+X^4$ term in the generating function numerator).
\item \textbf{Exact arithmetic casts and dual generating-function forms:} The closed form is stated over $\Q$ so division by~$2$ is exact field division; the recurrence casts $a(\cdot)$ to $\Z$ so $\N$-truncating subtraction cannot mask sign errors; and the generating function is pinned in $\Q\llbracket X \rrbracket$ in \emph{both} cleared-denominator polynomial form and explicit unit-inverse (\texttt{PowerSeries.inv}) form.
\item \textbf{Scope boundary (standalone formalization):} Although the ON-cell sibling \seq{A266256} (Seam~B entry \texttt{C052}) satisfies $a(n) + \seq{A266256}(n) = 2n + 1$, the pinned \seq{A266258} specification is proved completely standalone from its own verbatim definitions with zero cross-file sequence dependencies; companion Rule~47 OFF-cell sequence \seq{A266664} is discussed structurally in Remark~\ref{rem:rule47}.
\end{enumerate}

\subsection{Prior-Art Boundary}
\label{sec:prior-art}

Comprehensive literature and database audits conducted on 14~November~2026 (scoping lane: OEIS \seq{A266258} plus seven ECA siblings, Rule~11 cluster \seq{A266253}--\seq{A266259}, arXiv, OpenAlex, and web search) and re-verified live on 24~November~2026 (live Google web search via \texttt{search\_web} and Firecrawl SERPs, live OEIS text records, OpenAlex API, Zenodo API, Crossref API, and GitHub code search) confirm that Colin Barker's \seq{A266258} conjecture bundle remains unproved in the published mathematical literature and external formal proof libraries:
\begin{itemize}
\item On OEIS, \seq{A266258} remains at revision~\#14 (February~16, 2025) with Barker's \texttt{\%F} block explicitly labeled \emph{Conjectures from Colin Barker}, as do \seq{A266256} (\#15), \seq{A266259} (\#14), and \seq{A266664} (\#21). (On the row-encoding companions \seq{A266254} and \seq{A266255}, one-line formulas omit the word ``Conjectures'' but include zero proof text or citations.)
\item Sloane's survey on ON-cell counts in cellular automata~\cite{sloane2015} has zero occurrences of Rule~11 or \seq{A266258}.
\item Zero published papers or preprints citing \seq{A266258} appear on arXiv, OpenAlex, Zenodo, or Crossref. Public GitHub codebases mentioning \seq{A266258} ({\small\nolinkurl{archmageirvine/joeis}}, {\small\nolinkurl{loda-lang/loda-programs}}, {\small\nolinkurl{provables/sequencelib}}, and {\small\nolinkurl{ai4reason/oeis-atp-benchmark}}) were inspected first-hand and hold only a numerical Java CA simulator or synthetic term-fitting loop programs (tested up to $n \le 100$) with no formal Rule~11 definition and no proof linking Rule~11 evolution to Barker's formulas. Within \texttt{derekblair/math-proofs}, the sibling entry \texttt{C052} (\seq{A266256}, Rule~11 ON cells) proves the ON-cell bundle and the window complement law $a_{\mathrm{ON}}(n) + a_{\mathrm{OFF}}(n) = 2n + 1$, but does not state or prove Barker's \seq{A266258} closed form, recurrence, or rational generating function.
\end{itemize}

\subsection{Main Theorem}

\begin{definition}[Rule~11, bi-infinite $\Z$-evolution, and $a(n) = \seq{A266258}(n)$]
\label{def:c054}
Define $\ruleTT : \mathsf{Bool}^3 \to \mathsf{Bool}$ by the explicit $8$-branch truth table mapping $(\mathtt{false}, \mathtt{true}, \mathtt{true})$, $(\mathtt{false}, \mathtt{false}, \mathtt{true})$, and $(\mathtt{false}, \mathtt{false}, \mathtt{false})$ to $\mathtt{true}$ and the remaining five triples to $\mathtt{false}$. Define the global evolution $\ca : \N \to \Z \to \mathsf{Bool}$ recursively by
\[
\ca(0, j) := \mathrm{decide}(j = 0), \qquad
\ca(n + 1, j) := \ruleTT\bigl(\ca(n, j - 1),\, \ca(n, j),\, \ca(n, j + 1)\bigr).
\]
For every $n \in \N$, define the OFF-cell count on the active light-cone window $[-n, n] \subset \Z$ by
\[
a(n) := \bigl|\{j \in \Z \cap [-n, n] : \ca(n, j) = \mathtt{false}\}\bigr|.
\]
\end{definition}

\begin{theorem}[Main Theorem --- Resolution of Barker's \seq{A266258} Bundle]
\label{thm:main}
Let $\ruleTT$, $\ca$, and $a$ be as in Definition~\ref{def:c054}. Then:
\begin{enumerate}
\item[\textup{(1)}] \textbf{Initial value:} $a(0) = 0$.
\item[\textup{(2)}] \textbf{Rational closed form:} For every $n \in \N$ with $n > 0$,
\[
a(n) = \frac{2(-1)^n n + 2n - 3(-1)^n + 1}{2} \qquad \text{in }\Q.
\]
\item[\textup{(3)}] \textbf{Four-step linear recurrence:} For every $n \in \N$ with $n > 4$,
\[
a(n) = 2a(n - 2) - a(n - 4) \qquad \text{in }\Z.
\]
\item[\textup{(4)}] \textbf{Ordinary generating function (cleared-denominator and unit-inverse forms):} Let $F(X) := \sum_{n=0}^{\infty} a(n)\,X^n \in \Q\llbracket X \rrbracket$. Then in $\Q\llbracket X \rrbracket$,
\begin{align*}
(1 - X)^2 (1 + X)^2 \, F(X) &= X\bigl(2 + 3X - 2X^2 + X^3\bigr), \\
F(X) &= X\bigl(2 + 3X - 2X^2 + X^3\bigr) \cdot \bigl((1 - X)^2 (1 + X)^2\bigr)^{-1}.
\end{align*}
\end{enumerate}
Moreover, the guards $n > 0$ in \textup{(2)} and $n > 4$ in \textup{(3)} are both sharp.
\end{theorem}

\section{Persistent Two-Cell Glider and Global \texorpdfstring{$\Z$}{Z}-Invariant}

\begin{lemma}[Bit-index equivalence, Boolean normal form, and chirality]
\label{lem:rule11}
For all $(l, c, r) \in \mathsf{Bool}^3$:
\begin{enumerate}
\item[\textup{(a)}] $\ruleTT(l, c, r) = \mathrm{Nat.testBit}\;11\;(4l + 2c + r)$;
\item[\textup{(b)}] $\ruleTT(l, c, r) = (\neg l) \land (\neg c \lor r)$;
\item[\textup{(c)}] $\ruleTT$ breaks left-right reflection symmetry:
\[
\ruleTT(\mathtt{false}, \mathtt{false}, \mathtt{true}) = \mathtt{true} \ne \mathtt{false} = \ruleTT(\mathtt{true}, \mathtt{false}, \mathtt{false}).
\]
\end{enumerate}
\end{lemma}

\begin{proof}
Exhaustive evaluation over all $2^3 = 8$ Boolean triples $(l, c, r)$ verifies (a) and (b) directly against $11 = 00001011_2$: every triple with $l = \mathtt{true}$ outputs $\mathtt{false}$, and among the four triples with $l = \mathtt{false}$, only $(\mathtt{false}, \mathtt{true}, \mathtt{false})$ outputs $\mathtt{false}$, yielding $(\neg l) \land (\neg c \lor r)$. Part~(c) compares bit~$1$ ($1$) with bit~$4$ ($0$).
\end{proof}

\begin{proposition}[Generation $n = 1$ on $\Z$]
\label{prop:ca-one}
For every $j \in \Z$,
\[
\ca(1, j) = \mathrm{decide}(j \ne 0 \land j \ne 1).
\]
\end{proposition}

\begin{proof}
By Definition~\ref{def:c054} and Lemma~\ref{lem:rule11}(b), $\ca(1, j) = (\neg \mathrm{decide}(j - 1 = 0)) \land (\neg \mathrm{decide}(j = 0) \lor \mathrm{decide}(j + 1 = 0))$. Splitting into three exhaustive cases:
\begin{itemize}
\item $j = 0$: neighborhood $(\mathtt{false}, \mathtt{true}, \mathtt{false}) \mapsto \mathtt{false}$, matching $\mathrm{decide}(0 \ne 0 \land 0 \ne 1) = \mathtt{false}$.
\item $j = 1$: neighborhood $(\mathtt{true}, \mathtt{false}, \mathtt{false}) \mapsto \mathtt{false}$, matching $\mathrm{decide}(1 \ne 0 \land 1 \ne 1) = \mathtt{false}$.
\item $j \notin \{0, 1\}$: both $j - 1 \ne 0$ ($l = \mathtt{false}$) and $j \ne 0$ ($c = \mathtt{false}$), so $(\neg \mathtt{false}) \land (\neg \mathtt{false} \lor r) = \mathtt{true}$ regardless of $r$, matching $\mathrm{decide}(j \ne 0 \land j \ne 1) = \mathtt{true}$. \qedhere
\end{itemize}
\end{proof}

\begin{lemma}[Single-step glider transitions on $\Z$]
\label{lem:glider-steps}
Let $n \in \N$.
\begin{enumerate}
\item[\textup{(a)}] \emph{Odd-to-even transition:} If $\ca(n, j) = \mathrm{decide}(j \ne n - 1 \land j \ne n)$ for all $j \in \Z$, then
\[
\ca(n + 1, j) = \mathrm{decide}(j = n \lor j = n + 1) \qquad (\forall j \in \Z).
\]
\item[\textup{(b)}] \emph{Even-to-odd transition:} If $\ca(n, j) = \mathrm{decide}(j = n - 1 \lor j = n)$ for all $j \in \Z$, then
\[
\ca(n + 1, j) = \mathrm{decide}(j \ne n \land j \ne n + 1) \qquad (\forall j \in \Z).
\]
\end{enumerate}
\end{lemma}

\begin{proof}
Fix $j \in \Z$ and write $(l, c, r) = (\ca(n, j - 1), \ca(n, j), \ca(n, j + 1))$. By Lemma~\ref{lem:rule11}(b), $\ca(n + 1, j) = (\neg l) \land (\neg c \lor r)$.

\emph{Part~(a).} Under the hypothesis, $\ca(n, i) = \mathtt{false} \iff i \in \{n - 1, n\}$.
\begin{itemize}
\item If $j = n$, then $(j - 1, j, j + 1) = (n - 1, n, n + 1)$, so $(l, c, r) = (\mathtt{false}, \mathtt{false}, \mathtt{true}) \mapsto \mathtt{true}$.
\item If $j = n + 1$, then $(j - 1, j, j + 1) = (n, n + 1, n + 2)$, so $(l, c, r) = (\mathtt{false}, \mathtt{true}, \mathtt{true}) \mapsto \mathtt{true}$.
\item If $j \notin \{n, n + 1\}$, then $j - 1 \notin \{n - 1, n\}$, so $l = \mathtt{true}$ and $(\neg \mathtt{true}) \land (\neg c \lor r) = \mathtt{false}$.
\end{itemize}
Hence $\ca(n + 1, j) = \mathrm{decide}(j = n \lor j = n + 1)$.

\emph{Part~(b).} Under the hypothesis, $\ca(n, i) = \mathtt{true} \iff i \in \{n - 1, n\}$.
\begin{itemize}
\item If $j \in \{n, n + 1\}$, then $j - 1 \in \{n - 1, n\}$, so $l = \mathtt{true}$ and $\ca(n + 1, j) = \mathtt{false}$.
\item If $j = n - 1$, then $(j - 1, j, j + 1) = (n - 2, n - 1, n)$, so $(l, c, r) = (\mathtt{false}, \mathtt{true}, \mathtt{true}) \mapsto \mathtt{true}$.
\item If $j \notin \{n - 1, n, n + 1\}$, then both $j - 1 \notin \{n - 1, n\}$ ($l = \mathtt{false}$) and $j \notin \{n - 1, n\}$ ($c = \mathtt{false}$), so $(\neg \mathtt{false}) \land (\neg \mathtt{false} \lor r) = \mathtt{true}$ regardless of $r$.
\end{itemize}
Hence $\ca(n + 1, j) = \mathrm{decide}(j \ne n \land j \ne n + 1)$.
\end{proof}

\begin{theorem}[Global glider + alternating-exterior invariant on $\Z$]
\label{thm:global-inv}
For every $m \in \N$ and every $j \in \Z$:
\begin{align}
\ca(2m + 1, j) &= \mathrm{decide}(j \ne 2m \land j \ne 2m + 1), \label{eq:inv-odd}\\
\ca(2m + 2, j) &= \mathrm{decide}(j = 2m + 1 \lor j = 2m + 2). \label{eq:inv-even}
\end{align}
Equivalently, for every $n \in \N$ with $n \ge 1$ and every $j \in \Z$,
\begin{equation}
\label{eq:inv-xor}
\ca(n, j) = \operatorname{xor}\bigl(\mathrm{decide}(\mathrm{Odd}(n)),\, \mathrm{decide}(j = n - 1 \lor j = n)\bigr).
\end{equation}
\end{theorem}

\begin{proof}
Induct on $m \in \N$ for the coupled statement~\eqref{eq:inv-odd}--\eqref{eq:inv-even}.
\begin{itemize}
\item \emph{Base case $m = 0$:} Equation~\eqref{eq:inv-odd} at $m = 0$ is Proposition~\ref{prop:ca-one}, and Lemma~\ref{lem:glider-steps}(a) at $n = 1$ yields~\eqref{eq:inv-even} at $m = 0$.
\item \emph{Inductive step $m \mapsto m + 1$:} Assume~\eqref{eq:inv-even} holds at $m$, so $\ca(2m + 2, j) = \mathrm{decide}(j = (2m + 2) - 1 \lor j = 2m + 2)$ for all $j \in \Z$. Applying Lemma~\ref{lem:glider-steps}(b) at $n = 2m + 2$ yields~\eqref{eq:inv-odd} at $m + 1$, and applying Lemma~\ref{lem:glider-steps}(a) at $n = 2m + 3$ yields~\eqref{eq:inv-even} at $m + 1$.
\end{itemize}
For any $n \ge 1$, splitting on parity gives~\eqref{eq:inv-xor}: if $n = 2m$ ($m \ge 1$), $\mathrm{decide}(\mathrm{Odd}(n)) = \mathtt{false}$ and $\operatorname{xor}(\mathtt{false}, b) = b$, matching~\eqref{eq:inv-even} at $m - 1$; if $n = 2m + 1$, $\mathrm{decide}(\mathrm{Odd}(n)) = \mathtt{true}$ and $\operatorname{xor}(\mathtt{true}, \mathrm{decide}(j = 2m \lor j = 2m + 1)) = \mathrm{decide}(j \ne 2m \land j \ne 2m + 1)$, matching~\eqref{eq:inv-odd}.
\end{proof}

\section{OFF-Cell Counts, Closed Form, Recurrence, and Guard Sharpness}

\begin{proposition}[Exact OFF-cell parity counts and OFF/ON complement law]
\label{prop:counts}
For every $n \in \N$, the integer interval $[-n, n] \cap \Z$ has cardinality $2n + 1$. Moreover:
\begin{enumerate}
\item[\textup{(a)}] $a(0) = 0$.
\item[\textup{(b)}] For every $m \ge 0$, $a(2m + 1) = 2$.
\item[\textup{(c)}] For every $m \ge 1$, $a(2m) = 4m - 1 = 2(2m) - 1$.
\item[\textup{(d)}] Letting $a_{\mathrm{ON}}(n) := |\{j \in [-n, n] \cap \Z : \ca(n, j) = \mathtt{true}\}|$ (\seq{A266256}), we have $a(n) + a_{\mathrm{ON}}(n) = 2n + 1$ for every $n \in \N$.
\end{enumerate}
\end{proposition}

\begin{proof}
\begin{enumerate}
\item[\textup{(a)}] At $n = 0$, $[-0, 0] \cap \Z = \{0\}$ and $\ca(0, 0) = \mathtt{true} \ne \mathtt{false}$, so the filtered OFF set is $\emptyset$ of cardinality~$0$.
\item[\textup{(b)}] Let $m \ge 0$. By Theorem~\ref{thm:global-inv}, $\ca(2m + 1, j) = \mathtt{false} \iff j \in \{2m, 2m + 1\}$. Since $-(2m + 1) \le 2m < 2m + 1 \le 2m + 1$, the filtered finset equals the two-element set $\{2m, 2m + 1\}$ with $2m \ne 2m + 1$, giving $a(2m + 1) = 2$.
\item[\textup{(c)}] Let $m \ge 1$. By Theorem~\ref{thm:global-inv}, for $j \in \Z$,
\[
\bigl(-2m \le j \le 2m\bigr) \land \neg(j = 2m - 1 \lor j = 2m) \iff -2m \le j \le 2m - 2.
\]
Hence the filtered finset equals the contiguous integer interval $[-2m, 2m - 2] \cap \Z$, whose cardinality is $(2m - 2) + 1 - (-2m) = 4m - 1$.
\item[\textup{(d)}] Every $j \in [-n, n] \cap \Z$ satisfies either $\ca(n, j) = \mathtt{false}$ or $\ca(n, j) = \mathtt{true}$ exclusively, partitioning $[-n, n] \cap \Z$ into two disjoint subsets whose cardinalities sum to $|[-n, n] \cap \Z| = 2n + 1$. \qedhere
\end{enumerate}
\end{proof}

From Proposition~\ref{prop:counts} we immediately record the five initial values
\begin{equation}
\label{eq:init-vals}
a(0) = 0, \qquad a(1) = 2, \qquad a(2) = 3, \qquad a(3) = 2, \qquad a(4) = 7.
\end{equation}

\begin{theorem}[Rational closed form and sharpness at $n = 0$]
\label{thm:closed-form}
For every $n \in \N$ with $n > 0$,
\[
a(n) = \frac{2(-1)^n n + 2n - 3(-1)^n + 1}{2} \qquad \text{in }\Q.
\]
The guard $n > 0$ is sharp: at $n = 0$, $a(0) = 0 \ne -1 = \bigl(2(-1)^0 \cdot 0 + 2 \cdot 0 - 3(-1)^0 + 1\bigr)/2$.
\end{theorem}

\begin{proof}
Let $n > 0$ and split on parity.
\begin{itemize}
\item If $n = 2m$ ($m \ge 1$), then $(-1)^{2m} = 1$ and Proposition~\ref{prop:counts}(c) gives $a(2m) = 4m - 1$ (with $1 \le 4m$, so casting to $\Q$ preserves $4m - 1$). Thus
\[
\frac{2(1)(2m) + 2(2m) - 3(1) + 1}{2} = \frac{8m - 2}{2} = 4m - 1 = a(2m).
\]
\item If $n = 2m + 1$ ($m \ge 0$), then $(-1)^{2m+1} = -1$ and Proposition~\ref{prop:counts}(b) gives $a(2m + 1) = 2$. Thus
\[
\frac{2(-1)(2m + 1) + 2(2m + 1) - 3(-1) + 1}{2} = \frac{4}{2} = 2 = a(2m + 1).
\]
\end{itemize}
Sharpness at $n = 0$ follows directly from $a(0) = 0 \ne (-3 + 1)/2 = -1$.
\end{proof}

\begin{theorem}[Four-step linear recurrence and sharpness at $n = 4$]
\label{thm:recurrence}
For every $n \in \N$ with $n > 4$,
\[
a(n) = 2a(n - 2) - a(n - 4) \qquad \text{in }\Z.
\]
The guard $n > 4$ is sharp: at $n = 4$, $a(4) = 7 \ne 6 = 2a(2) - a(0)$.
\end{theorem}

\begin{proof}
Since $n > 4$, we have $n \ge 5$, so $n - 2 \ge 3$ and $n - 4 \ge 1$ are strictly positive and share the parity of~$n$:
\begin{itemize}
\item If $n = 2m$ ($m \ge 3$), then $n - 2 = 2(m - 1)$ and $n - 4 = 2(m - 2)$ with $m, m - 1, m - 2 \ge 1$. Proposition~\ref{prop:counts}(c) gives
\[
2a(n - 2) - a(n - 4) = 2\bigl(4(m - 1) - 1\bigr) - \bigl(4(m - 2) - 1\bigr) = 4m - 1 = a(n).
\]
\item If $n = 2m + 1$ ($m \ge 2$), then $n - 2 = 2(m - 1) + 1$ and $n - 4 = 2(m - 2) + 1$. Proposition~\ref{prop:counts}(b) gives $a(n) = a(n - 2) = a(n - 4) = 2$, whence $2(2) - 2 = 2 = a(n)$.
\end{itemize}
At $n = 4$,~\eqref{eq:init-vals} gives $a(4) = 7$ whereas $2a(2) - a(0) = 2(3) - 0 = 6$.
\end{proof}

\section{Rational Ordinary Generating Function in \texorpdfstring{$\Q\llbracket X \rrbracket$}{Q[[X]]}}

\begin{theorem}[Cleared-denominator generating function]
\label{thm:gf}
Let $F(X) := \sum_{n=0}^{\infty} a(n)\,X^n \in \Q\llbracket X \rrbracket$. Then
\begin{equation}
\label{eq:gf-cleared}
(1 - X)^2 (1 + X)^2 \, F(X) = X\bigl(2 + 3X - 2X^2 + X^3\bigr) \qquad \text{in }\Q\llbracket X \rrbracket.
\end{equation}
\end{theorem}

\begin{proof}
In the commutative ring $\Q\llbracket X \rrbracket$, polynomial expansion yields
\begin{align*}
(1 - X)^2 (1 + X)^2 \, F(X) &= (1 - 2X^2 + X^4)\,F(X) = F(X) - 2(X^2 F(X)) + X^4 F(X), \\
X\bigl(2 + 3X - 2X^2 + X^3\bigr) &= 2X + 3X^2 - 2X^3 + X^4.
\end{align*}
By extensionality of formal power series, it suffices to check equality of the $n$-th coefficient $[X^n]$ for every $n \in \N$. Using $[X^n](X^d F(X)) = a(n - d)$ if $n \ge d$ and $0$ otherwise, we verify the six cases $n = 0, 1, 2, 3, 4$ and $n = k + 5$ ($k \ge 0$) directly via~\eqref{eq:init-vals} and Theorem~\ref{thm:recurrence}:
\begin{align*}
n = 0:\quad & a(0) = 0 = [X^0](2X + 3X^2 - 2X^3 + X^4), \\
n = 1:\quad & a(1) = 2 = [X^1](2X + 3X^2 - 2X^3 + X^4), \\
n = 2:\quad & a(2) - 2a(0) = 3 - 2(0) = 3 = [X^2](2X + 3X^2 - 2X^3 + X^4), \\
n = 3:\quad & a(3) - 2a(1) = 2 - 2(2) = -2 = [X^3](2X + 3X^2 - 2X^3 + X^4), \\
n = 4:\quad & a(4) - 2a(2) + a(0) = 7 - 2(3) + 0 = 1 = [X^4](2X + 3X^2 - 2X^3 + X^4), \\
n = k + 5:\quad & a(k + 5) - 2a(k + 3) + a(k + 1) = 0 = [X^{k+5}](2X + 3X^2 - 2X^3 + X^4),
\end{align*}
where the last line follows by casting Theorem~\ref{thm:recurrence} at $n = k + 5 > 4$ from $\Z$ to $\Q$.
\end{proof}

\begin{corollary}[Unit invertibility and explicit inverse-series generating function]
\label{cor:unit-inv}
The polynomial factor $D(X) := (1 - X)^2(1 + X)^2 \in \Q\llbracket X \rrbracket$ is a unit in $\Q\llbracket X \rrbracket$, and
\begin{equation}
\label{eq:gf-inv}
F(X) = X\bigl(2 + 3X - 2X^2 + X^3\bigr) \cdot \bigl((1 - X)^2 (1 + X)^2\bigr)^{-1} \qquad \text{in }\Q\llbracket X \rrbracket.
\end{equation}
\end{corollary}

\begin{proof}
A formal power series over a field is a unit if and only if its constant coefficient is nonzero. Here $[X^0]D(X) = (1 - 0)^2(1 + 0)^2 = 1 \ne 0$, so $D(X) \cdot D(X)^{-1} = 1$ in $\Q\llbracket X \rrbracket$. Multiplying~\eqref{eq:gf-cleared} on the right by $D(X)^{-1}$ yields $F(X) = (D(X) F(X))\,D(X)^{-1} = X(2 + 3X - 2X^2 + X^3)\,D(X)^{-1}$.
\end{proof}

\begin{proof}[Proof of Theorem~\ref{thm:main}]
Combine Proposition~\ref{prop:counts}(a), Theorem~\ref{thm:closed-form}, Theorem~\ref{thm:recurrence}, Theorem~\ref{thm:gf}, and Corollary~\ref{cor:unit-inv}.
\end{proof}

\section{Structural Extensions: Rule~43 Equivalence, Rule~47 Contrast, and Zero-Clamping Traps}

\subsection{Isolated-Zero Avoidance and the \texorpdfstring{$\{11, 43\}$}{\{11, 43\}} Equivalence Class}

For any Wolfram rule code $c \in \N$, let $\operatorname{rule}_c(l, c_0, r) := \mathrm{Nat.testBit}\;c\;(4l + 2c_0 + r)$ and let $\ca_c : \N \to \Z \to \mathsf{Bool}$ denote its bi-infinite single-seed evolution on $\Z$.

\begin{theorem}[Isolated-zero avoidance and Rule~11 / Rule~43 equivalence]
\label{thm:rule43}
\begin{enumerate}
\item[\textup{(a)}] For every $n \in \N$ and every $j \in \Z$, the isolated-zero neighborhood $(\mathtt{true}, \mathtt{false}, \mathtt{true})$ (bit index~$5$) never occurs in $\ca(n, \cdot)$:
\[
\neg\bigl(\ca(n, j - 1) = \mathtt{true} \land \ca(n, j) = \mathtt{false} \land \ca(n, j + 1) = \mathtt{true}\bigr).
\]
\item[\textup{(b)}] If $c \in \N$ agrees with~$11$ on all seven bits $i \in \{0, 1, 2, 3, 4, 6, 7\}$, then $\ca_c(n, j) = \ca(n, j)$ for all $n \in \N$ and $j \in \Z$. In particular, both Wolfram Rule~$11$ ($00001011_2$) and Wolfram Rule~$43$ ($00101011_2$) generate bit-for-bit identical space-time diagrams on $\Z$ and identical OFF-cell counts $a(n) = \seq{A266258}(n)$.
\end{enumerate}
\end{theorem}

\begin{proof}
For (a), inspect the three regimes of $\ca(n, \cdot)$:
\begin{itemize}
\item $n = 0$: $\ca(0, j - 1) = \mathtt{true} = \ca(0, j + 1)$ requires $j - 1 = 0 = j + 1$, impossible.
\item $n = 2m \ge 2$ even: by Theorem~\ref{thm:global-inv}, $\ca(2m, \cdot)$ is $\mathtt{true}$ solely on the adjacent pair $\{2m - 1, 2m\}$, so $j - 1, j + 1 \in \{2m - 1, 2m\}$ cannot hold simultaneously since $(j + 1) - (j - 1) = 2 > 1$.
\item $n = 2m + 1 \ge 1$ odd: by Theorem~\ref{thm:global-inv}, $\ca(2m + 1, \cdot)$ is $\mathtt{false}$ solely on $\{2m, 2m + 1\}$. If $\ca(2m + 1, j) = \mathtt{false}$, then $j \in \{2m, 2m + 1\}$, so either $j + 1 = 2m + 1$ ($\ca(n, j + 1) = \mathtt{false}$) or $j - 1 = 2m$ ($\ca(n, j - 1) = \mathtt{false}$).
\end{itemize}
Part~(b) follows by induction on $n \in \N$: at each step $n \mapsto n + 1$, the induction hypothesis gives $\ca_c(n, \cdot) = \ca(n, \cdot)$, and by part~(a) bit~$5$ is never queried, so $\operatorname{rule}_c$ agrees with $\ruleTT$. Since $43 = 11 + 2^5$, this formally establishes the $\{11, 43\}$ equivalence noted on \seq{A266253}~\cite{oeis266253}.
\end{proof}

\begin{remark}[Contrast with defect-annihilating rules and Rule~47 (\seq{A266664})]
\label{rem:rule47}
Unlike Rule~7 (\seq{A266220}/\seq{A266222}) and Rule~23 (\seq{A266437}/\seq{A266439}), which annihilate the seed defect by $n = 2$ and $n = 1$ respectively to yield spatially uniform alternating rows on $\Z$, Rule~11 sustains the persistent two-cell glider $\{n - 1, n\}$ for all $n \ge 1$. Meanwhile, Wolfram Rule~47 ($47 = 00101111_2$, OFF-cell count \seq{A266664}~\cite{oeis266664}) differs from Rule~11 at bit~$2$ ($010 \mapsto 1$) and bit~$5$ ($101 \mapsto 1$). At $n = 1$, bit~$2$ flips $j = 0$ from $0$ to $1$, yielding window \texttt{110} (a single OFF cell at $j = 1$, so $\seq{A266664}(1) = 1 \ne 2 = a(1)$). At $n = 2$, however, $j = 1$ sees the isolated-zero triple $(1, 0, 1)$, which bit~$5$ of Rule~47 maps to~$1$, while $j = 2$ sees $(0, 1, 1) \mapsto 1$ and all other cells map to~$0$---producing the exact same two-cell glider $\{1, 2\}$ (\texttt{00011}) as Rule~11. Thereafter bit~$5$ is never queried again, explaining both R.~J.\ Mathar's identity $\seq{A266664}(n) = \seq{A266258}(n)$ for $n > 1$ and why Rule~47's sharp guards shift up by~$+1$ ($n > 1$ and $n > 5$) relative to Rule~11 ($n > 0$ and $n > 4$).
\end{remark}

\subsection{Causal-Cone Mirror and Formalization of Both Zero-Clamping Traps}

Direct unmemoized evaluation of $\ca(n, j)$ via Definition~\ref{def:c054} traverses a ternary recursion tree of size $3^n$. To enable fast, pure-kernel spot-checking of $a(n)$ directly against OEIS~\seq{A266258} without any computational trust escape, we define an $O(n^2)$ causal-cone contraction mirror $a_{\mathrm{dp}}(n)$: initialize the length-$(4n + 1)$ Boolean list $[0^{2n}, 1, 0^{2n}]$ representing $\ca(0, j)$ on $[-2n, 2n]$, iterate the sliding-window map $(l, c, r, \dots) \mapsto (\ruleTT(l, c, r), \dots)$ exactly $n$ times (shrinking length by~$2$ per step down to $2n + 1$), and count the $\mathtt{false}$ entries.

\begin{proposition}[Causal-cone mirror and both zero-clamping traps]
\label{prop:clamping-traps}
\begin{enumerate}
\item[\textup{(a)}] For every $n \in \N$, $a_{\mathrm{dp}}(n) = a(n)$, and pure kernel reduction confirms
\[
\bigl(a(0), a(1), \dots, a(12)\bigr) = (0, 2, 3, 2, 7, 2, 11, 2, 15, 2, 19, 2, 23),
\]
matching OEIS~\seq{A266258}.
\item[\textup{(b)}] \emph{Trap~A (light-cone zero-clamping):} Let $\ca_{\mathrm{cone}}$ evolve by Rule~11 while clamping every cell outside $[-m, m]$ to $\mathtt{false}$ after step~$m$, and let $a_{\mathrm{cone}}(n)$ count its OFF cells on $[-n, n]$. Then $a_{\mathrm{cone}}(2) = 2 \ne 3 = a(2)$.
\item[\textup{(c)}] \emph{Trap~B (fixed-box $[-n, n]$ zero-clamping):} Let $\ca_{\mathrm{box}}(R, \cdot, \cdot)$ evolve by Rule~11 on the fixed box $[-R, R]$ with exterior clamped to $\mathtt{false}$, and let $a_{\mathrm{box}}(n)$ count OFF cells of $\ca_{\mathrm{box}}(n, n, \cdot)$ on $[-n, n]$. Then $\ca_{\mathrm{box}}(2, 2, -2) = \mathtt{true} \ne \mathtt{false} = \ca(2, -2)$ (bit profile already corrupted at $n = 2$: \texttt{10010} vs.\ \texttt{00011}) and $a_{\mathrm{box}}(3) = 4 \ne 2 = a(3)$.
\end{enumerate}
\end{proposition}

\begin{proof}
Part~(a) follows by induction on the step count $k$: applying $k$ sliding-window steps to $[\ca(m, s), \dots, \ca(m, s + \ell + 2k - 1)]$ yields $[\ca(m + k, s + k), \dots, \ca(m + k, s + k + \ell - 1)]$; setting $m = 0$, $s = -2n$, $\ell = 2n + 1$, and $k = n$ recovers the exact window $[-n, n]$ at generation $n$, and filtering for $\mathtt{false}$ matches $a(n)$. Parts~(b) and~(c) follow by direct kernel evaluation: under Trap~A, row~$1$ is \texttt{100} with zero exterior, producing row~$2$ equal to \texttt{10011} (OFF count~$2 \ne 3$); under Trap~B at $R = 2$, row~$1$ on $[-2, 2]$ is \texttt{11001} with zero exterior, producing row~$2$ equal to \texttt{10010} ($\ca_{\mathrm{box}}(2, 2, -2) = \mathtt{true}$), and at $R = 3$, row~$3$ on $[-3, 3]$ is \texttt{0011100} (OFF count~$4 \ne 2$).
\end{proof}

\section{Lean 4 Formalization and Verification Record}

The formalization resides in the repository \texttt{derekblair/math-proofs} under Lean~\texttt{v4.34.1} (\nolinkurl{leanprover/lean4:v4.34.1}) and Mathlib revision {\small\nolinkurl{d13f23b723b8a846827a245b89c10fc7d3f11612}}. Following the live Seam~B repository convention, the files reconcile the pre-Seam-B target path \nolinkurl{Mathproofs/ECARule11OffC054.lean} from \texttt{PIN-C054.md} while preserving the pinned namespace \nolinkurl{Mathproofs.ECARule11OffC054} and every definition and theorem signature verbatim:
\begin{itemize}\setlength{\itemsep}{2pt}
\item \nolinkurl{MathProofs/SeamB/C054/Statement.lean} (107 lines): verbatim transcription of frozen \texttt{PIN-C054.md} v1 (\texttt{delta\_vs\_source: NONE}).\\
Statement SHA-256: {\small\nolinkurl{a6971e3dea0fed9853a9c7c3cab5884cff286cfbd0066d2a84f19a92dc29f171}}.\\
Pin SHA-256: {\small\nolinkurl{93a607c738a4c5d08bb5c139179fadc8b8ef04a68061eac8599b1283d85fa5a0}}.
\item \nolinkurl{MathProofs/SeamB/C054/Solution.lean} (773 lines): reproduces \texttt{rule11}, \texttt{ca}, \texttt{a}, and all six pinned theorem signatures verbatim in \nolinkurl{Mathproofs.ECARule11OffC054}, proving every pinned, structural, and non-vacuity theorem with zero \texttt{sorry}s, zero warnings, and allowlist-only axioms.\\
Solution SHA-256: {\small\nolinkurl{8c921a639a828c1f24bbcb26751ce60808de8e69ce283cbe8ca233f73af7acce}}.
\item \nolinkurl{Mathproofs.lean} (97 lines; adds one import for \nolinkurl{MathProofs.SeamB.C054.Solution}).\\
Root SHA-256: {\small\nolinkurl{3b412e935a189b63927c66e1d7a10a2e9151cdbc5e43ae89c2782e4af9395027}}.
\end{itemize}

Table~\ref{tab:lean-map} records the correspondence between the mathematical results in this paper and the machine-checked declarations in \nolinkurl{MathProofs/SeamB/C054/Solution.lean}.

\begin{table}[ht]
\centering
\caption{Correspondence between paper results and Lean~4 declarations in \nolinkurl{Mathproofs.ECARule11OffC054}.}
\label{tab:lean-map}
\smallskip
{\small
\begin{tabular}{@{}l >{\raggedright\arraybackslash}p{8.2cm} l@{}}
\toprule
Paper Result & Lean~4 Declaration(s) & Status \\
\midrule
Definition~\ref{def:c054} & \texttt{rule11}, \texttt{ca}, \texttt{a} (pinned verbatim) & Exported \\
Lemma~\ref{lem:rule11} & \texttt{rule11\_eq\_testBit}, \texttt{rule11\_eq}, \texttt{rule11\_not\_symm} & Exported, proved \\
Proposition~\ref{prop:ca-one} & \texttt{ca\_zero}, \texttt{ca\_one} & Exported, proved \\
Lemma~\ref{lem:glider-steps} & \texttt{ca\_step\_odd\_to\_even}, \texttt{ca\_step\_even\_to\_odd} & Exported, proved \\
Theorem~\ref{thm:global-inv} & \texttt{ca\_odd\_even}, \texttt{ca\_odd}, \texttt{ca\_even\_pos}, \texttt{ca\_eq\_xor} & Exported, proved \\
Proposition~\ref{prop:counts} & \texttt{a266258\_zero} (pinned), \texttt{a\_odd}, \texttt{a\_even\_pos}, \texttt{a\_add\_aOn} & Exported, proved \\
Theorem~\ref{thm:closed-form} & \texttt{a266258\_closed\_form} (pinned, guard $0 < n$), \texttt{a266258\_closed\_form\_false\_zero} & Exported, proved \\
Theorem~\ref{thm:recurrence} & \texttt{a266258\_recurrence} (pinned, guard $4 < n$), \texttt{a266258\_recurrence\_false\_four} & Exported, proved \\
Theorem~\ref{thm:gf} & \texttt{a266258\_gf} (pinned) & Exported, proved \\
Corollary~\ref{cor:unit-inv} & \texttt{a266258\_gf\_denom\_isUnit}, \texttt{a266258\_gf\_inv} (pinned) & Exported, proved \\
Theorem~\ref{thm:main} & \texttt{c054\_a266258\_bundle} (pinned) & Exported, proved \\
Theorem~\ref{thm:rule43} & \texttt{ca\_no\_101}, \texttt{caRule\_eq\_ca}, \texttt{caRule\_equiv\_class} & Exported, proved \\
Proposition~\ref{prop:clamping-traps} & \texttt{a\_dp\_eq\_a}, \texttt{a\_dp\_spot\_0\_to\_12}, \texttt{a\_spot\_0\_to\_12}, \texttt{aConeClamp\_two\_eq\_two}, \texttt{aBoxClamp\_trap} & Exported, proved \\
\bottomrule
\end{tabular}
}
\end{table}

\noindent\textbf{Audit summary.}
\begin{enumerate}\setlength{\itemsep}{2pt}
\item \emph{Clean-slate build:} Running \texttt{rm -rf .lake/build \&\& lake build} on \nolinkurl{MathProofs.SeamB.C054.Solution} exits with code~$0$, zero warnings, and 1{,}795 jobs ($2.3\text{ s}$ module build, $3.72\text{ s}$ wall).
\item \emph{Independent kernel replay:} Running \texttt{lake env leanchecker -{}-fresh -{}-verbose} on the compiled module \nolinkurl{MathProofs.SeamB.C054.Solution} exits with code~$0$ ($153.70\text{ s}$ wall).
\item \emph{Axiom audit:} \texttt{\#print axioms} on every pinned theorem and a \texttt{CoreM} \texttt{collectAxioms} sweep across all 224 declarations in {\small\nolinkurl{Mathproofs.ECARule11OffC054}} confirm zero disallowed axioms ({\small\texttt{DISALLOWED\_AXIOMS = []}}): every declaration depends strictly on subsets of \texttt{[propext, Classical.choice, Quot.sound]} (in particular, \texttt{rule11}, \texttt{ca}, {\small\texttt{rule11\_eq}}, {\small\texttt{rule11\_not\_symm}}, and {\small\texttt{ca\_zero}} depend on zero axioms \texttt{[]}; {\small\texttt{rule11\_eq\_testBit}} and {\small\texttt{a\_dp\_spot\_0\_to\_12}} depend solely on \texttt{[propext]}).
\item \emph{Statement fidelity:} Automated diff confirms 9/9 verbatim definition and signature matches across \texttt{PIN-C054.md} v1, \texttt{Statement.lean}, and \texttt{Solution.lean}, and kernel \texttt{rfl} definitional-equation plus direct term-assignment type-identity verification (exit~$0$) confirms exact identity against \texttt{PIN-C054.md} v1. (The external \texttt{comparator} tool targeting Lean \texttt{v4.36.0} was not run because this repository pins \texttt{v4.34.1} and \texttt{comparator} is not installed on the build host; the gap is recorded explicitly in \texttt{VERIFICATION.md}.)
\item \emph{Lexical trust-escape scan:} $0$ raw hits and $0$ comment/string-stripped hits for \texttt{sorry}, \texttt{admit}, \texttt{sorryAx}, \texttt{axiom}, \texttt{native\_decide}, \texttt{ofReduceBool}, \texttt{ofReduceNat}, \texttt{implemented\_by}, and \texttt{unsafe} in \texttt{Solution.lean} and \texttt{Mathproofs.lean}.
\end{enumerate}

{\small
\begin{thebibliography}{99}\setlength{\itemsep}{2pt}
\raggedright

\bibitem{oeis266258}
R.\ Price, \emph{OEIS sequence \seq{A266258}}: Number of OFF (white) cells in the $n$-th iteration of the ``Rule~11'' elementary cellular automaton starting with a single ON (black) cell. Includes conjectures by C.\ Barker (December~27, 2015 and April~14, 2019); live revision \#14 (February~16, 2025). \url{https://oeis.org/A266258}.

\bibitem{oeis266253}
R.\ Price, \emph{OEIS sequence \seq{A266253}}: Triangle read by rows giving successive states of cellular automaton generated by ``Rule~11'' initiated by a single ON (black) cell. Live revision \#17 (August~8, 2025). \url{https://oeis.org/A266253}.

\bibitem{oeis266254}
R.\ Price, \emph{OEIS sequence \seq{A266254}}: Binary representations of the $n$-th iteration of the ``Rule~11'' elementary cellular automaton starting with a single ON (black) cell. Live revision \#24 (February~16, 2025). \url{https://oeis.org/A266254}.

\bibitem{oeis266255}
R.\ Price, \emph{OEIS sequence \seq{A266255}}: Decimal representations of the $n$-th iteration of the ``Rule~11'' elementary cellular automaton starting with a single ON (black) cell. Live revision \#22 (February~16, 2025). \url{https://oeis.org/A266255}.

\bibitem{oeis266256}
R.\ Price, \emph{OEIS sequence \seq{A266256}}: Number of ON (black) cells in the $n$-th iteration of the ``Rule~11'' elementary cellular automaton starting with a single ON (black) cell. Includes conjectures by C.\ Barker (December~27, 2015 and April~14, 2019); live revision \#15 (February~16, 2025). \url{https://oeis.org/A266256}.

\bibitem{oeis266664}
R.\ Price, \emph{OEIS sequence \seq{A266664}}: Number of OFF (white) cells in the $n$-th iteration of the ``Rule~47'' elementary cellular automaton starting with a single ON (black) cell. Includes comment by R.~J.\ Mathar (January~10, 2016); live revision \#21 (February~16, 2025). \url{https://oeis.org/A266664}.

\bibitem{wolfram2002}
S.\ Wolfram, \emph{A New Kind of Science}, Wolfram Media, Champaign, IL, 2002 (Section~2, p.~55).

\bibitem{mathworldECA}
E.~W.\ Weisstein, \emph{Elementary Cellular Automaton}, from MathWorld---A Wolfram Web Resource. \url{https://mathworld.wolfram.com/ElementaryCellularAutomaton.html}.

\bibitem{sloane2015}
N.~J.~A.\ Sloane, \emph{On the Number of ON Cells in Cellular Automata}, arXiv:1503.01168 (2015). \url{https://arxiv.org/abs/1503.01168}.

\end{thebibliography}
}

\end{document}
